\chapter{Loads and Boundary Conditions}
\label{chap:loads-bcs}

Loads and boundary conditions act on the compiled Assembly. Native FEMaster separates reusable load/support definitions from the analysis that activates them: \cmd{SUPPORT} entries are collected in named support collectors and native load commands are collected in named load collectors. A \cmd{LOADCASE} later selects one or more collectors with \cmd{SUPPORTS} and \cmd{LOADS}. This makes the same physical loading reusable across several analysis definitions.

Abaqus input follows Step semantics instead. \cmd{BOUNDARY}, \cmd{CLOAD}, \cmd{DLOAD}, and \cmd{DSLOAD} inside a supported \cmd{STEP} are translated into FEMaster's internal support/load objects and inserted into parser-owned Step collectors. The analysis-procedure chapter explains how those Step collectors are attached to the resulting FEMaster loadcase.

\keywordpage{AMPLITUDE}

\cmd{AMPLITUDE} defines a reusable scalar function of analysis time. Native FEMaster supports \val{TYPE=LINEAR}, \val{STEP}, and \val{NEAREST}. Every data row contains one time/value pair. Linear interpolation is the default; step interpolation holds the preceding sample, while nearest interpolation selects the closest sample according to the amplitude implementation.

The supported Abaqus reader maps \cmd{AMPLITUDE, DEFINITION=TABULAR} to the same internal amplitude object using linear interpolation. It currently accepts \key{TIME=STEPTIME} and \key{VALUE=RELATIVE}. Abaqus data lines may contain up to four time/value pairs per physical line (eight scalars); each nonempty pair becomes a sample. Other Abaqus amplitude definitions and time bases are outside the registered subset.

An amplitude does not apply a load by itself. A load references it by name. Dynamic procedures can retain the amplitude for evaluation at physical time, whereas supported static Abaqus translations may sample it according to the procedure semantics implemented by the reader.

\begin{keywordparameters}
\key{NAME} & required & Unique amplitude name. \\
\key{TYPE} & native optional & \val{LINEAR}, \val{STEP}, or \val{NEAREST}; default \val{LINEAR}. \\
\key{DEFINITION} & Abaqus optional & Supported value \val{TABULAR}; default \val{TABULAR}. \\
\key{TIME} & Abaqus optional & Supported value \val{STEPTIME}. \\
\key{VALUE} & Abaqus optional & Supported value \val{RELATIVE}. \\
\end{keywordparameters}

\begin{keyworddata}
Native & One \val{time, value} pair per record. \\
Abaqus & Up to four \val{time, value} pairs per physical line. \\
\end{keyworddata}

\exampleheading

\begin{dialectexamples}
\begin{femasterdeck}
*AMPLITUDE, NAME=RAMP, TYPE=LINEAR
0.0, 0.0
0.1, 1.0
1.0, 1.0
\end{femasterdeck}
\begin{abaqusdeck}
*AMPLITUDE, NAME=RAMP,
 DEFINITION=TABULAR, TIME=STEPTIME
0.0, 0.0, 0.1, 1.0,
1.0, 1.0
\end{abaqusdeck}
\end{dialectexamples}

\notesheading

Amplitude time is dimensionally the same time variable used by the active analysis. For transient calculations this is physical integration time; for static mappings the Abaqus reader follows its documented Step-time scaling rather than inventing a dynamic history.

\keywordpage{SUPPORT}

Native \cmd{SUPPORT} defines prescribed generalized nodal components and stores them in a named support collector. Each row targets either a compiled node set or one resolvable node reference and supplies six entries in the order \dofvec. A finite value prescribes that component; an omitted entry or \val{NAN} remains unconstrained. Zero therefore means a real homogeneous constraint and must not be used when the intended component is free.

\key{SUPPORT_COLLECTOR} selects the collector. An optional \key{ORIENTATION} resolves a named coordinate system and causes the six prescribed components to be interpreted in that local basis. The support object retains the region, generalized values, and orientation; constraint equations are generated later when a loadcase activates the collector.

The Abaqus analogue is \cmd{BOUNDARY}, documented separately below. Abaqus rows use one-based DOF numbers rather than six simultaneous values, but both forms are translated into the same six-component FEMaster support representation.

\begin{keywordparameters}
\key{SUPPORT_COLLECTOR} & required & Named native support collector. \\
\key{ORIENTATION} & optional & Named local coordinate system for the prescribed components. \\
\end{keywordparameters}

\begin{keyworddata}
\val{target} & Node set or resolvable node reference. \\
Six DOF values & \(U_x,U_y,U_z,R_x,R_y,R_z\); missing/empty values become NaN (free). \\
\end{keyworddata}

\exampleheading

\begin{dialectexamples}
\begin{femasterdeck}
*SUPPORT, SUPPORT_COLLECTOR=BC
FIXED, 0., 0., 0., 0., 0., 0.
GUIDE, NAN, 0., 0., NAN, NAN, NAN
\end{femasterdeck}
\begin{abaqusdeck}
*BOUNDARY
FIXED, 1, 6, 0.
GUIDE, 2, 3, 0.
\end{abaqusdeck}
\end{dialectexamples}

\notesheading

A target node set is not copied when the support is parsed; the support retains a shared region reference. This permits a collector to contain several independently oriented supports while the loadcase handles them uniformly through the chosen constraint method.

\keywordpage{BOUNDARY}

\cmd{BOUNDARY} is the supported Abaqus displacement/rotation constraint command. It may appear at root level for persistent zero constraints or inside a \cmd{STEP} after a supported procedure has activated the FEMaster loadcase. The optional \key{TYPE} currently accepts only \val{DISPLACEMENT} and defaults to that form.

Each row contains a target, the first structural DOF, an optional last DOF, and an optional magnitude. DOFs are one based and must lie in the range 1--6. If the last DOF is omitted, only the first DOF is prescribed. If magnitude is omitted it defaults to zero. FEMaster converts the selected range into a six-component support vector with NaN in all unspecified positions.

If an Abaqus \cmd{TRANSFORM} was assigned to a target node, the generated support uses that local nodal coordinate system. For a node set the reader creates the appropriately oriented support contribution for each compiled member.

Nonzero prescribed values inside a Step are currently supported only for static procedures. An amplitude on a nonzero prescribed value is further restricted to supported linear static translation, where the reader evaluates the amplitude according to the Step period. Amplitudes outside a Step are not supported by the current translation.

\begin{keywordparameters}
\key{TYPE} & optional & Supported form \val{DISPLACEMENT}; default \val{DISPLACEMENT}. \\
\key{AMPLITUDE} & optional & Supported under the procedure restrictions described above. \\
\end{keywordparameters}

\begin{keyworddata}
\val{target, first_dof, last_dof, magnitude} & Structural DOF range 1--6; last DOF defaults to first and magnitude defaults to zero. \\
\end{keyworddata}

\exampleheading

\begin{dialectexamples}
\begin{femasterdeck}
*SUPPORT, SUPPORT_COLLECTOR=BC
TIP, 1.5, 0., 0., NAN, NAN, NAN
\end{femasterdeck}
\begin{abaqusdeck}
*BOUNDARY, TYPE=DISPLACEMENT
TIP, 1, 1, 1.5
TIP, 2, 3, 0.
\end{abaqusdeck}
\end{dialectexamples}

\notesheading

The right-hand example creates two Abaqus boundary rows because the Abaqus grammar addresses contiguous DOF ranges while native FEMaster presents the full generalized vector in one record.

\keywordpage{CLOAD}

A concentrated load applies generalized nodal force and moment components. Native \cmd{CLOAD} stores one six-component vector per target in a named load collector. The order is \(F_x,F_y,F_z,M_x,M_y,M_z\), corresponding to translational and rotational nodal DOFs. Missing entries default to zero.

The target may be a node set or one resolvable node reference. \key{ORIENTATION} optionally gives a local coordinate system and \key{AMPLITUDE} optionally attaches a named time function. Transformation and time scaling are intentionally deferred to load evaluation rather than baked into the stored vector.

Abaqus \cmd{CLOAD} uses one target, one DOF number, and one magnitude per row. FEMaster translates that row into a six-component vector with one nonzero entry. Abaqus structural DOFs 1--6 are supported. Nodal \cmd{TRANSFORM} assignments determine the local load basis. \key{AMPLITUDE} is resolved according to the active procedure. \key{FOLLOWER} is not supported, imaginary loads are rejected, and a CLOAD is not permitted in a FREQUENCY eigenproblem Step.

\begin{keywordparameters}
\key{LOAD_COLLECTOR} & native required & Destination native load collector. \\
\key{ORIENTATION} & native optional & Named local load basis. \\
\key{AMPLITUDE} & optional & Named scalar amplitude where supported. \\
\key{FOLLOWER} & Abaqus flag & Currently rejected. \\
\key{REAL}/\key{IMAGINARY} & Abaqus flags & Real loading is supported; imaginary loading is not. \\
\end{keywordparameters}

\exampleheading

\begin{dialectexamples}
\begin{femasterdeck}
*CLOAD, LOAD_COLLECTOR=TIP_LOAD
TIP, 1000., -250., 0., 0., 0., 0.
\end{femasterdeck}
\begin{abaqusdeck}
*CLOAD
TIP, 1, 1000.
TIP, 2, -250.
\end{abaqusdeck}
\end{dialectexamples}

\notesheading

When a target is a node set, the prescribed vector is applied to each node in that set; it is not automatically divided by the number of nodes. Use a distributing coupling when the modelling intent is to distribute a resultant load according to coupling weights.

\keywordpage{DLOAD}

The name \cmd{DLOAD} has different surface syntax in the two dialects and must therefore be read in context. Native FEMaster \cmd{DLOAD} is a vector traction on a surface or surface set. Each target receives a three-component traction vector; optional \key{ORIENTATION} and \key{AMPLITUDE} references are retained. Surface integration and consistent nodal-force reduction occur during assembly.

The supported Abaqus \cmd{DLOAD} mapping currently implements only the \val{GRAV} body-acceleration form. Its data contain an element/element-set target, the label \val{GRAV}, a scalar magnitude, and a three-component direction. FEMaster multiplies magnitude and direction (together with any procedure-dependent amplitude scale) and creates its internal volume-load object. An empty target resolves to \val{EALL}. Other Abaqus distributed-load labels and imaginary loading are rejected.

For surface tractions and pressure in Abaqus input, use the supported \cmd{DSLOAD} command documented below rather than interpreting native FEMaster's \cmd{DLOAD} semantics as Abaqus semantics.

\begin{keywordparameters}
\key{LOAD_COLLECTOR} & native required & Destination native load collector. \\
\key{ORIENTATION} & native optional & Local traction basis. \\
\key{AMPLITUDE} & optional & Named time scaling. \\
\key{REAL}/\key{IMAGINARY} & Abaqus flags & Imaginary distributed loading is not supported. \\
\end{keywordparameters}

\exampleheading

\begin{dialectexamples}
\begin{femasterdeck}
*DLOAD, LOAD_COLLECTOR=SHEAR
SIDE, 0., 2.5, 0.
\end{femasterdeck}
\begin{abaqusdeck}
** Abaqus surface vector traction uses DSLOAD.
*DSLOAD, FOLLOWER=NO
SIDE, TRVEC, 2.5, 0., 1., 0.
\end{abaqusdeck}
\end{dialectexamples}

\keywordpage{PLOAD}

Native \cmd{PLOAD} applies scalar pressure to a surface or surface set. Pressure direction follows the finite-element surface normal; the command therefore needs no orientation parameter. The pressure magnitude may reference a named amplitude. Surface normals, Jacobians, and consistent nodal forces are evaluated when the load is assembled.

Abaqus input expresses the supported equivalent through \cmd{DSLOAD} with load type \val{P}. The Abaqus row contains the named surface, \val{P}, and the scalar pressure. Direction components are forbidden for the pressure form. The current Abaqus translation rejects follower pressure in nonlinear static Steps, because that combination is not represented by the implemented nonlinear load formulation.

\begin{keywordparameters}
\key{LOAD_COLLECTOR} & native required & Destination load collector. \\
\key{AMPLITUDE} & optional & Named pressure scaling function. \\
\end{keywordparameters}

\exampleheading

\begin{dialectexamples}
\begin{femasterdeck}
*PLOAD, LOAD_COLLECTOR=PRESSURE
INNER_FACE, 2.0
\end{femasterdeck}
\begin{abaqusdeck}
*DSLOAD
INNER_FACE, P, 2.0
\end{abaqusdeck}
\end{dialectexamples}

\keywordpage{DSLOAD}

\cmd{DSLOAD} is the supported Abaqus surface-load command. Two load labels are implemented: \val{P} for scalar pressure and \val{TRVEC} for a vector traction. The target must be a compiled named surface set.

For \val{P}, the data contain surface name, load type, and magnitude; no direction components may be supplied. FEMaster creates the same internal pressure load used by native \cmd{PLOAD}. For \val{TRVEC}, three finite direction components are required. The direction vector is normalized and multiplied by the specified magnitude, then stored as FEMaster's internal vector surface traction. An optional \key{ORIENTATION} can supply the local traction basis.

\key{FOLLOWER} defaults to \val{YES}. In nonlinear static Steps, \val{TRVEC} is currently supported only with \key{FOLLOWER=NO}; follower pressure is likewise not supported in nonlinear Steps. Real loading is supported and imaginary loading is rejected. \key{AMPLITUDE} follows the active-procedure amplitude resolution.

\begin{keywordparameters}
\key{AMPLITUDE} & optional & Named amplitude. \\
\key{ORIENTATION} & optional & Named basis for \val{TRVEC}. \\
\key{FOLLOWER} & optional & \val{YES} or \val{NO}; default \val{YES}, subject to nonlinear restrictions. \\
\key{REAL}/\key{IMAGINARY} & optional flags & Imaginary loading is unsupported. \\
\end{keywordparameters}

\begin{keyworddata}
Pressure & \val{surface, P, magnitude}. \\
Traction & \val{surface, TRVEC, magnitude, dx, dy, dz}. \\
\end{keyworddata}

\exampleheading

\begin{dialectexamples}
\begin{femasterdeck}
*PLOAD, LOAD_COLLECTOR=LOADS
FACE_A, 5.0
*DLOAD, LOAD_COLLECTOR=LOADS
FACE_B, 0., 3.0, 0.
\end{femasterdeck}
\begin{abaqusdeck}
*DSLOAD, FOLLOWER=NO
FACE_A, P, 5.0
FACE_B, TRVEC, 3.0, 0., 1., 0.
\end{abaqusdeck}
\end{dialectexamples}

\keywordpage{VLOAD}

Native \cmd{VLOAD} applies a three-component distributed volume/body-force vector to an element or element set. The load is integrated over element volume and reduced to generalized nodal forces by the element/load implementation. Optional orientation and amplitude references follow the same deferred-evaluation semantics as native surface tractions.

The supported Abaqus analogue is \cmd{DLOAD} with the \val{GRAV} label. The Abaqus translation multiplies scalar acceleration magnitude by the supplied direction and creates the same internal FEMaster volume-load class. Thus the native vector is the already resolved acceleration/body-force vector, while the Abaqus syntax stores magnitude and direction separately.

\begin{keywordparameters}
\key{LOAD_COLLECTOR} & native required & Destination load collector. \\
\key{ORIENTATION} & native optional & Named local vector basis. \\
\key{AMPLITUDE} & optional & Named time scaling. \\
\end{keywordparameters}

\exampleheading

\begin{dialectexamples}
\begin{femasterdeck}
*VLOAD, LOAD_COLLECTOR=GRAVITY
EALL, 0., 0., -9810.
\end{femasterdeck}
\begin{abaqusdeck}
*DLOAD
EALL, GRAV, 9810., 0., 0., -1.
\end{abaqusdeck}
\end{dialectexamples}

\keywordpage{TLOAD}

Native \cmd{TLOAD} creates a thermal load from an already compiled scalar nodal temperature Field and a stress-free reference temperature. Unlike mechanical load commands, all information is supplied through keyword parameters: \key{LOAD_COLLECTOR}, \key{TEMPERATUREFIELD}, and \key{REFERENCETEMPERATURE}.

The referenced Field must use the \val{NODE} domain and exactly one component. During analysis, elements combine the local temperature difference with their material thermal-expansion coefficient and constitutive response. \cmd{TLOAD} itself does not precompute thermal strains or equivalent forces because those depend on element geometry and material properties.

The supported Abaqus reader currently has no generic temperature/thermal-load translation corresponding to this native mechanism.

\begin{keywordparameters}
\key{LOAD_COLLECTOR} & required & Destination load collector. \\
\key{TEMPERATUREFIELD} & required & Existing scalar \val{NODE} Field. \\
\key{REFERENCETEMPERATURE} & required & Stress-free reference temperature. \\
\end{keywordparameters}

\exampleheading

\begin{dialectexamples}
\begin{femasterdeck}
*FIELD, NAME=TEMPERATURE,
 TYPE=NODE, COLS=1, FILL=ZERO
1, 80.
2, 100.
*TLOAD, LOAD_COLLECTOR=THERMAL,
 TEMPERATUREFIELD=TEMPERATURE,
 REFERENCETEMPERATURE=20.
\end{femasterdeck}
\begin{noabaqus}
No corresponding thermal-load command is currently registered by the supported Abaqus reader.
\end{noabaqus}
\end{dialectexamples}

\keywordpage{INERTIALOAD}

The native command is spelled \cmd{INERTIALOAD} in the current registry. It defines rigid-body inertial acceleration contributions over an element set. Each data row contains the target element set followed by four three-component vectors: reference center, translational acceleration of that center, angular velocity \(\boldsymbol\omega\), and angular acceleration \(\boldsymbol\alpha\).

The load implementation evaluates the rigid-body acceleration field and converts the mass distribution into consistent inertial forces during assembly. This can represent translational inertia, centrifugal contributions due to angular velocity, and tangential contributions due to angular acceleration. \key{CONSIDER_POINT_MASSES} controls whether native concentrated point-mass features participate; its default is false.

The target must be an existing compiled element set. Mass density and any included point masses must be physically consistent for the inertial load to have meaningful magnitude. The supported Abaqus reader currently has no direct translation for this full native command.

\begin{keywordparameters}
\key{LOAD_COLLECTOR} & required & Destination load collector. \\
\key{CONSIDER_POINT_MASSES} & optional & Boolean; default false/0. \\
\end{keywordparameters}

\begin{keyworddata}
\val{target} & Existing element set. \\
\val{cx,cy,cz} & Rigid-body reference center. \\
\val{ax,ay,az} & Translational acceleration of the center. \\
\val{wx,wy,wz} & Angular velocity. \\
\val{alphax,alphay,alphaz} & Angular acceleration. \\
\end{keyworddata}

\exampleheading

\begin{dialectexamples}
\begin{femasterdeck}
*INERTIALOAD, LOAD_COLLECTOR=ROT_INERTIA,
 CONSIDER_POINT_MASSES=1
EALL,
0., 0., 0.,
0., 0., 0.,
0., 0., 10.,
0., 0., 2.
\end{femasterdeck}
\begin{noabaqus}
No direct equivalent is currently registered by the supported Abaqus reader.
\end{noabaqus}
\end{dialectexamples}

\notesheading

The spelling \cmd{INERTIALOAD} contains a single ``L'' between inertial and load because that is the exact command currently registered by FEMaster. The keyword index uses the executable grammar as the source of truth and therefore preserves this spelling.
